\section{Types in plai-typed}
\label{sec:types}

\code{plai-typed} is installed with \code{raco pkg install plai-typed}, and a file starts with \code{\#lang plai-typed}. It looks like Racket, but a type checker reads the whole program first, and a program that fails the check does not run at all.

\subsection{Every Value Has a Type}

At the REPL, every answer comes with two lines: the type the checker worked out, starting with \code{-}, and then the value.
\begin{lstlisting}[style=repl]
> 3
- number
3
> (+ 3 4)
- number
7
> "hello"
- string
"hello"
> (+ 3 "hello")
typecheck failed: number vs. string
\end{lstlisting}
This is the idea of the lecture in miniature: the type is known \emph{before} the value is computed. The checker does not run \rkt|(+ 3 4)| to find out it is a number; it knows because \code{+} takes numbers and returns a number.

\subsection{Type Inference and Annotations}

No types need to be written. Here the checker sees \rkt|(+ 3 x)|, concludes that \code{x} must be a number, and that the result is one too:
\begin{lstlisting}
(define (add3 x)
  (+ 3 x))
\end{lstlisting}

\begin{definition}[Type inference]
\emph{Type inference} is the checker working out the types that the programmer did not write down.
\end{definition}

Types can also be written, as many or as few as we like. An argument becomes \rkt|(x : number)|, and the return type goes after the argument list. Whatever is left out is inferred:
\begin{lstlisting}
(define (add3-v2 (x : number))            ; return type inferred
  (+ 3 x))

(define (add3-v3 (x : number)) : number   ; both annotated
  (+ 3 x))
\end{lstlisting}

\begin{pitfall}[Spaces around the colon]
\code{(x: number)} is read as one identifier named \code{x:}, followed by \code{number}, and is rejected with \code{expected either <id> or `[<id> : <type>]'}. Write \code{(x : number)}.
\end{pitfall}

An annotation is a claim the checker verifies, not a hint it trusts. This one is wrong, since \rkt|(+ 3 x)| is a number, so the definition is rejected:
\begin{lstlisting}
(define (add3-v4 (x : number)) : string
  (+ 3 x))
;; typecheck failed: number vs. string
\end{lstlisting}
Either way the program is checked before it runs: a call like \rkt|(add3-v3 "hello")| never gets the chance to fail at run time, because the whole file is rejected.

If the checker can infer the types, why write them? Two reasons from lecture:
\begin{itemize}
  \item \textbf{Documentation.} The reader sees what the function takes and returns without reading its body.
  \item \textbf{Mental model.} Writing the type first makes you decide what the function is supposed to do before you write how it does it.
\end{itemize}

\subsection{Function Types}
\label{sec:fntypes}

Functions are values, so they have types. Read the arrow as ``takes \ldots{} and returns \ldots'':
\begin{lstlisting}[style=repl]
> add3
- (number -> number)
\end{lstlisting}
With two arguments, both go to the left of the arrow. Annotations can be mixed: in \code{myadd-v2}, \code{y} is unannotated but added to a number, so it is inferred to be one, and the two functions have the same type.
\begin{lstlisting}
(define (myadd x y)
  (+ x y))
(define (myadd-v2 (x : number) y) : number
  (+ x y))
\end{lstlisting}
\begin{lstlisting}[style=repl]
> myadd
- (number number -> number)
> myadd-v2
- (number number -> number)
\end{lstlisting}

A function that returns a function, \code{make-adder} from \Cref{sec:hof}, now has a type that says so:
\begin{lstlisting}
(define (myadd-v3 (x : number))
  (lambda (y) (+ x y)))
\end{lstlisting}
\begin{lstlisting}[style=repl]
> myadd-v3
- (number -> (number -> number))
> (myadd-v3 2)
- (number -> number)
> ((myadd-v3 2) 5)
- number
7
\end{lstlisting}
Compare the two types. \code{myadd} takes two numbers at once. \code{myadd-v3} takes \emph{one} number and returns a function that waits for the second, so \rkt|(myadd-v3 2)| is not a number yet but a function.

\subsection{Lists Are Homogeneous}

A list type says what its elements are, and every element must have that type:
\begin{lstlisting}[style=repl]
> (list 1 2 3 4)
- (listof number)
'(1 2 3 4)
> (cons 3 (cons 4 empty))
- (listof number)
'(3 4)
> (list 1 "hello")
typecheck failed: number vs. string
> (cons 3 4)
typecheck failed: (listof number) vs. number
\end{lstlisting}
Plain Racket accepts both of the last two. In \code{plai-typed} the second argument of \code{cons} must be a list, so bare pairs (\Cref{sec:lists}) are impossible. In exchange, list operations have precise types: the checker knows \rkt|(first (list 1 2 3 4))| is a \code{number} and \rkt|(rest (list 1 2 3 4))| a \code{(listof number)}, and likewise for \code{second} and \code{third}, without running anything.

\subsection{Pairs}

When two values of different types belong together, use a \emph{pair} (a tuple). Its type is written with \code{*}:
\begin{lstlisting}[style=repl]
> (pair 1 "hello")
- (number * string)
'#(1 "hello")
> (fst (pair 1 2))
- number
1
> (snd (pair 1 2))
- number
2
> (car (pair 1 2))
car: unbound identifier
\end{lstlisting}
A list has any length but one element type; a pair has exactly two slots, each with its own type. A pair is not a cons cell, and \code{plai-typed} provides no \code{car} or \code{cdr} at all: use \code{fst}/\code{snd} for pairs and \code{first}/\code{rest} for lists.

\subsection{Polymorphic Types}

Asking for the type of \code{first} shows something new:
\begin{lstlisting}[style=repl]
> first
- ((listof 'a) -> 'a)
> rest
- ((listof 'a) -> (listof 'a))
> cons
- ('a (listof 'a) -> (listof 'a))
> fst
- (('a * 'b) -> 'a)
> map
- (('a -> 'b) (listof 'a) -> (listof 'b))
\end{lstlisting}

\begin{definition}[Type variable, parametric polymorphism]
\code{'a} (read ``alpha'') is a \emph{type variable}: any type may stand in its place, but every \code{'a} in one type must stand for the \emph{same} type. A single definition that works uniformly at every type in this way is \emph{parametrically polymorphic}; other languages call it generics.
\end{definition}

So \code{first} takes a \code{(listof number)} to a \code{number} and a \code{(listof string)} to a \code{string}, but can never take a list of numbers to a string. Different letters allow different types: \code{fst} takes a pair of an \code{'a} and a \code{'b} and returns the \code{'a}. And the type of \code{map} alone says most of what \code{map} does: given a function from \code{'a} to \code{'b} and a list of \code{'a}, it produces a list of \code{'b}.

\subsection{Why Insist on Homogeneous Lists?}

Suppose \rkt|(list 1 "hello")| were allowed; call it \code{l}. What is the type of \rkt|(first l)|? It depends on which element happens to come first, and the checker cannot know that before the program runs. Then
\begin{itemize}
  \item every use of \code{first} would need a check at run time, which is exactly what a type checker is supposed to spare us, and
  \item \rkt|(map f l)| is only safe if \code{f} accepts \emph{both} numbers and strings, which \code{('a -> 'b)} cannot express.
\end{itemize}
One could design a type that lists every position, something like \code{(listof number string number)} for \rkt|(list 1 "hello" 2)|, but then every list function (\code{map}, \code{filter}, the folds) would need a far more complicated type. The simpler fix keeps lists homogeneous and instead makes \emph{one} type that covers all the kinds of things to be mixed. That is what \code{define-type} is for.
